\documentclass[10pt,a4paper]{amsart}
\usepackage[margin=19mm]{geometry}
\usepackage{amsmath,amssymb,mathtools}
\usepackage[scr=rsfs,cal=euler]{mathalfa}
\usepackage{xfrac}
\usepackage[unicode,hidelinks,bookmarks=false]{hyperref}
\newcommand{\ie}{i.e.\ }
\newcommand{\cf}{cf.\ }
\newcommand{\resp}{respectively\ }
\newcommand{\Subsection}{\subsection}
\providecommand{\nobreakdash}{\nobreak\hbox{-}}
\setlength{\emergencystretch}{2em}
\multlinegap0pt
\usepackage{multirow}
\usepackage{xcolor}
\usepackage{amssymb}
\usepackage{booktabs}
\usepackage{algorithm}\usepackage{algpseudocode}
\usepackage[shortlabels]{enumitem}
\usepackage{tikz}
\newtheorem{theorem}{Theorem}
\title{On some F-spaces associated to the functions satisfying Mulholland inequality.}
\author{Lav Kumar Singh, Aljosa Peperko}
\date{Source: arXiv:2507.01661v2 (CC BY 4.0)}
\begin{document}
\maketitle

\noindent\emph{Source and licence.} On some F-spaces associated to the functions satisfying Mulholland inequality.. Version arXiv:2507.01661v2 (CC BY 4.0), distributed by arXiv under the Creative Commons Attribution 4.0 International licence, http://creativecommons.org/licenses/by/4.0/, as recorded in the arXiv licence metadata for this exact version and shown on the abstract page https://arxiv.org/abs/2507.01661v2. Full licence notice: \texttt{licenses/arxiv-2507.01661-CC-BY-4.0.txt}.\par\smallskip

\noindent\textit{Editorial scope.} Counted unit: the article's theorem locb (source0.tex lines 382--384). The article's complete original proof of it is delivered with it (source0.tex lines 385--408, one \texttt{proof} environment of 24 lines). No mathematics is written, repaired or re-derived: the statement and the proof are the article's own lines, in the article's own notation, and no retained line is substituted or re-flowed.  No further source-local context is needed: every cross-reference of every retained line resolves inside the two delivered spans.  Nothing was omitted from any retained span, so the package carries no omission record.  No retained line contains a citation, so the wrapper carries no bibliography and no bibliography entry.  No retained line contains a figure environment, an image-inclusion command or drawing code, so no image is delivered and the package declares no asset dependency.  Editorial formatting only, in addition: an AMS \texttt{amsart} preamble that restates the article's own declarations and macros used by the retained lines, namely the environments \texttt{theorem} on separate counters, together with the standard packages those lines need.  The retained environments print in this document as Theorem 1 in order of appearance, whereas the article prints the same items with the numbering of its own full document.  The complete native source of the article as distributed in the arXiv e-print for this exact version is bundled unchanged as \texttt{sources/180895/source0.tex}.
	 \begin{theorem}\label{locb}
	 	For a family of locally bounded $F$-spaces $\{X_i\}_{i=1}^n$ and an Young function $\Phi$ which satisfies Mulholland condition, the $F$-space $\left(\oplus_nX_i,d_\Phi\right)$ is locally bounded.
	 \end{theorem}

	 \begin{proof}
	 	It would be sufficient to prove this for the case $n=2$. Let $B_r^{(i)}$ denote the open ball of the F-space $X_i$ of radius $r$ and centered at $0$.  The open ball of $\left(X\oplus X_2,d_\Phi\right)$ centered at $(0,0)$ and of radius $r$ will be denoted by $B_r^\Phi$. Notice that for $(x,y)\in B_1$, we have 
	 	$\Phi^{-1}\left(\Phi(\|x\|)+\Phi(\|y\|)\right)< 1$. Thus $\Phi(\|x\|)<\Phi(1)$ and $\Phi(\|y\|)<1$. Hence, $\|x\|<1$ and similarly $\|y\|<1$. Thus, $(x,y)\in B_1^{(1)}\times B_1^{(2)}$ and 
	 	\begin{equation}\label{inc}
	 		B_1^\Phi\subset B_1^{(1)}\times B_1^{(2)}.
	 	\end{equation}
	 	Further, if $c,r>0$ then for any $(cx,cy)\in cB_{r}^{(1)}\times cB_{r}^{(2)}$, we have 
	 	\begin{align}
	 		\left|\left|\left(\frac{cx}{2c}.\frac{cy}{2c}\right)\right|\right|_\Phi&=\Phi^{-1}\left(\Phi(\frac{\|cx\|}{2c})+\Phi(\frac{\|cy\|}{2c})\right)\nonumber\\&\leq\Phi^{-1}\left(\frac{1}{2}\Phi(\|x\|)+\frac{1}{2}\Phi(\|y\|)\right)&&\text{due to convexity of}~\Phi\nonumber\\&\leq \Phi^{-1}\left(\frac{1}{2}\Phi(r)+\frac{1}{2}\Phi(r)\right)&&\text{due to}~\Phi~\text{being increasing}\nonumber\\&=r.
	 	\end{align}
	 	Thus, \begin{equation}\label{inc1}
cB_{r}^{(1)}\times cB_{r}^{2}\subset 2cB_r^{\Phi}~~~~~~~~~~~~~~~~~~~~~~~~~~~~~~~~~~~~~\text{for each}~c,r>0
	 	\end{equation}
	 	Now fix a $r_o<1$. Since $X_1$ and $X_2$ are locally, we can find a $t$ such that $B_1^{(1)}\subset sB_{r_o}^{(1)}$ and $B_1^{(2)}\subset sB_{r_o}^{(2)}$ for all $s>t$. Hence, 
	 	\begin{equation}
	 		B_1^{(1)}\times B_1^{(2)}\subset sB_{r_o}^{(1)}\times sB_{r_o}^{(2)}~~~~~~~~~~~~~~\forall s>t.
	 	\end{equation}
	 	
	 	Combining the above inclusion with the inclusion in equation-\ref{inc} and \ref{inc1}, we get 
	 	\begin{equation}
B_1^\Phi\subset	B_1^{(1)}\times B_1^{(2)}\subset sB_{r_o}^{(1)}\times sB_{r_o}^{(2)}\subset 2sB_{r_o}^\Phi~~\forall s>t.
	 	\end{equation}
	 	Thus, the collection of open balls of radius less than one and centered at origin forms a neighborhood base at origin. The set $B_1$ is bounded and hence $\left(X_1\oplus X_2,d_\Phi\right)$ is a locally bounded $F$-space.
	 \end{proof}

\end{document}
